% ============================================================
%  Section 4.x — Implementation Details and Evaluation Protocol
%  Goes at the END of Section 4 (Experimental Setup), after the
%  partitioning subsection. Holds every hyperparameter and metric
%  definition, so that Section 3 stays generic.
%  Values verified against src/gan.py, src/model.py,
%  src/multitask.py, src/conformal.py and notebooks 04 / 08.
% ============================================================

\subsection{Implementation Details and Evaluation Protocol}
\label{sec:implementation}

The framework of Section~\ref{sec:methodology} is instantiated on
XJTU-SY with the values listed in Table~\ref{tab:hyperparams}. The
staging parameters of Eqs.~\eqref{eq:fpt_threshold}--\eqref{eq:fpt}
are those motivated in Section~\ref{sec:hi_staging}: the baseline is
estimated over the first $\eta = 20\%$ of each bearing's life, the
relative floor is $\rho = 0.20$, and the persistence length is $m = 5$
recordings. Windows contain $w = 32$ recordings of $d = 20$ features.
The augmentation ratio $r = 0.75$ and the stage-loss weight
$\lambda = 1.0$ were selected on preliminary runs over
$r \in \{0.25, 0.5, 0.75, 1.0\}$ and $\lambda \in \{0.3, 0.5, 1.0\}$.

\begin{table}[htbp]
\centering
\caption{Hyperparameters of the framework components.}
\label{tab:hyperparams}
\footnotesize
\renewcommand{\arraystretch}{1.15}
\begin{tabular}{lll}
\toprule
\textbf{Component} & \textbf{Setting} & \textbf{Value} \\
\midrule
\multirow{8}{*}{C1 (generator, critic)}
 & Noise dimension & 64 \\
 & Conv.\ channels / kernel & 128 / 5 \\
 & BiLSTM units per direction & 128 \\
 & Optimiser, learning rate & Adam, $10^{-4}$ \\
 & Adam $(\beta_1, \beta_2)$ & $(0.5, 0.9)$ \\
 & Batch size / epochs & 64 / 300 \\
 & Critic steps $n_{\mathrm{c}}$ & 3 \\
 & Penalty weight $\lambda_{\mathrm{gp}}$ & 10 \\
\midrule
\multirow{8}{*}{C2 (predictor)}
 & BiLSTM layers / units & 2 / 48 \\
 & Head hidden units & 48 \\
 & Dropout & 0.4 \\
 & Optimiser, learning rate & Adam, $5 \times 10^{-4}$ \\
 & Weight decay & $5 \times 10^{-4}$ \\
 & Batch / max.\ epochs / patience & 64 / 100 / 25 \\
 & Gradient clipping (norm) & 5.0 \\
 & Stage-loss weight $\lambda$ & 1.0 \\
\midrule
Augmentation & Ratio $r$ & 0.75 \\
C3 (conformal) & Miscoverage $\alpha$ & 0.10\textsuperscript{a} \\
\bottomrule
\end{tabular}

\vspace{2pt}
\raggedright\scriptsize \textsuperscript{a}Headline level; 0.05, 0.15
and 0.20 are also evaluated.
\end{table}

All models were implemented in PyTorch and trained on a single NVIDIA
T4 GPU. Recurrent training on the GPU is not bit-wise reproducible,
and the RUL error of a given configuration was observed to vary by
about $\pm 0.015$--$0.02$ in normalised units from one run to the
next. Every configuration is therefore trained under five random seeds
in each of the five folds, and results are reported as the mean and
standard deviation over seeds. A difference smaller than this
seed-to-seed variability is not interpreted as an effect.\\

Point accuracy is measured by the root mean square error on the
normalised RUL scale. It is computed for each test bearing and then
averaged over bearings, so that the long-lived bearings, which supply
most of the windows, do not dominate the result:
\begin{equation}
  \mathrm{RMSE} = \frac{1}{|\mathcal{B}_{\mathrm{te}}|}
  \sum_{b \in \mathcal{B}_{\mathrm{te}}}
  \sqrt{ \frac{1}{|\mathcal{W}_b|} \sum_{t \in \mathcal{W}_b}
  \left( \hat{y}^{b}_{t} - y^{b}_{t} \right)^2 } ,
  \label{eq:rmse}
\end{equation}
where $\mathcal{W}_b$ is the set of evaluated windows of bearing $b$.
Interval quality is measured by the prediction interval coverage
probability and the mean prediction interval width,
\begin{equation}
  \mathrm{PICP} = \frac{1}{n} \sum_{i=1}^{n}
  \mathbb{1}\left\{ y_i \in \hat{C}(\mathbf{X}_i) \right\},
  \qquad
  \mathrm{MPIW} = \frac{1}{n} \sum_{i=1}^{n}
  \left| \hat{C}(\mathbf{X}_i) \right| ,
  \label{eq:picp_mpiw}
\end{equation}
computed over the test windows of each fold and averaged over folds.
PICP should lie at or slightly above the nominal level $1-\alpha$, and
a smaller MPIW is preferable only when PICP is maintained; the two
are therefore always reported together. All three metrics are
reported separately for the healthy, early-degradation and
near-failure stages, and for the post-FPT and complete trajectories.
The near-failure stage is the primary subset, as it is the regime that
motivates the framework and determines maintenance decisions.
